\section{本章小结}
\label{sec:minimax-summary}

本章研究固定代价$f$下的问题$\inf_{\bm{x}}\sup_{\bm{y}}f$，它有两层结构：外层最小化，内层取最坏情形。开头的三个困难，分别由最坏值函数、鞍点和更新动力学来回答，这三者回答的是不同的问题。

\begin{BookInlineTable}
 \centering
 \begin{tabular}{@{}lll@{}}
  \toprule
  工具 & 回答的问题 & 依赖的条件\\
  \midrule
  最坏值函数$\phi$ & 外层在最小化什么 & $\bm{Y}$紧时连续，响应唯一时可微\\
  全局鞍点与间隙 & 一对候选是否全局解 & 紧凸且凸凹时存在\\
  迭代动力学 & 算法会不会靠近解 & 单调性、Lipschitz性、步长\\
  \bottomrule
 \end{tabular}
 \caption{本章的三个主要工具。}
 \label{tab:minimax-summary}
\end{BookInlineTable}

顺序极小极大问题与同时下降—上升算法是两回事。前者是要解的数学问题，后者是一种迭代格式。图~\ref{fig:minimax-saddle-types}(c)的原点是顺序问题的解，却不是鞍点；$f=xy$上的同时更新发散，问题本身却有鞍点。

判断一个结果可信，要分别核对四项条件：
\begin{enumerate}
 \item 最值交换。只有凸凹性和相应的紧性成立时才能交换$\inf$和$\sup$，值相等还要另外核对达到，才有鞍点。
 \item 内层准确性。Danskin公式要全局最优响应，局部极大不能代替。强凹时内层间隙能转成外层梯度的误差，非凹时梯度残差小不保证接近最坏值。
 \item 耦合迭代的稳定。反对称的耦合产生旋转，下降—上升在双线性问题上不收敛。内层强凹时要让$\bm{y}$跟得上，单调问题上可以用外梯度或乐观梯度校正方向。
 \item 最终误差。间隙是全局证书，但只用近似内层算出的间隙偏小，需要修正；限定间隙和局部二阶条件都不能冒充全局结论。
\end{enumerate}

本章两层共用同一个代价$f$，内层想让$f$大，外层想让$f$小。下一章的两层目标各不相同，外层依赖内层的最优解，那时值函数的导数不再由这样简单的公式给出。
